\documentclass[11pt,a4paper]{ctexart}

% ── Packages ──
\usepackage{amsmath}
\usepackage{amssymb}
\usepackage{booktabs}
\usepackage{listings}
\usepackage{xcolor}
\usepackage[hidelinks]{hyperref}
\usepackage{graphicx}
\usepackage{authblk}

% ── Metadata ──
\title{\bfseries Whisper：面向Token效率的AI原生编程语言}
\author[1]{Myoucai}
\affil[1]{独立研究者}
\date{2026年6月}

\begin{document}

\maketitle

\begin{abstract}
大语言模型（LLM）越来越多地被用于代码生成，但其经济可行性受限于按Token计价的推理成本。
现有的通用语言（如Python）带有大量语法开销，显著增加了Token消耗。
本文提出\textbf{Whisper}——一门从零开始为最小化Token使用而设计的栈式后缀编程语言，
同时保持了表达能力和安全性。\textbf{Whisper}采用单字符操作符实现常用栈操作
（\texttt{\_} 复制、\texttt{`} 交换、\texttt{@} 旋转），紧凑的条件语法
（\texttt{?? ... | ... ]}），以及基于能力的零信任安全模型。
我们在1,149条\textbf{Whisper}指令-响应对上微调了Qwen2.5-Coder-7B模型。
在15个任务的代码生成基准测试上，微调后的模型实现了40.0\%的精确匹配准确率
和93.3\%的语法有效性。在另外25个任务的Token效率基准测试上，生成的
\textbf{Whisper}代码比等效Python代码少消耗58.8\%的Token。
在五个真实编程场景中，\textbf{Whisper}在4/5的任务上
Token效率优于Python，平均减少37.8\%。实验结果表明，面向LLM Token经济性的
语言级优化可以在不牺牲正确性的前提下，带来显著的成本节约。
\end{abstract}

% ═══════════════════════════════════════════════════════════════
\section{引言}
% ═══════════════════════════════════════════════════════════════

AI辅助软件开发生态的经济基础是以Token为单位的计费模型。每一次函数调用、
每一次代码补全、每一次模块生成，其成本都与底层语言模型处理的Token数量成正比。
在大规模场景下——企业每天数百万次API调用——即使10\%的Token消耗减少也能
转化为可观的财务节省。

然而，LLM最常生成的编程语言——Python、JavaScript、Java——都是为人类可读性
而非机器生成效率而设计的。它们的冗长带来了隐性成本：每个关键字
（\texttt{def}、\texttt{function}、\texttt{class}）、每对分隔符
（\texttt{( )}、\texttt{\{ \}}）、每个命名参数都消耗Token，这些Token贡献了
推理成本，却未增加超出模型已推断出的语义价值。

我们引入\textbf{Whisper}——一门从零设计为LLM代码生成\textit{原生输出格式}
的编程语言。\textbf{Whisper}具备以下特点：

\begin{itemize}
\item \textbf{Token最小化}：每个语法结构都经过选择以最小化Token数量。常用操作
  使用单字符操作符。整个语法可以用不到30个不同的Token来表达。
\item \textbf{栈式后缀表示}：受Forth和PostScript启发，\textbf{Whisper}使用
  栈评估模型和后缀表示法。这消除了括号、优先级规则、变量声明以及其他
  消耗Token但不增加信息的语法结构。
\item \textbf{默认安全}：除非通过能力Token系统显式授权，否则\textbf{Whisper}
  程序没有任何I/O能力。这使其适用于在不受信任的环境中执行AI生成的代码。
\item \textbf{自举}：\textbf{Whisper}编译器本身用\textbf{Whisper}编写，
  证明了该语言的实际表达能力。
\end{itemize}

本文做出以下贡献：

\begin{enumerate}
\item \textbf{Whisper}的设计与实现——一门面向LLM Token效率优化的新型栈式语言
  （第\ref{sec:design}节）。
\item 基于能力的零信任安全模型，支持AI生成代码的安全执行
  （第\ref{sec:security}节）。
\item 实证评估表明，微调后的7B参数模型在\textbf{Whisper}代码生成上实现了
  40.0\%的精确匹配准确率和93.3\%的语法有效性（第\ref{sec:experiments}节）。
\item 跨25个编程任务的Token效率基准测试，显示\textbf{Whisper}代码比等效Python
  少消耗58.8\%的Token（第\ref{sec:results}节）。
\end{enumerate}

% ═══════════════════════════════════════════════════════════════
\section{语言设计}
\label{sec:design}
% ═══════════════════════════════════════════════════════════════

\textbf{Whisper}是一门拼接式、基于栈的后缀表示语言。每个程序都是操作隐式数据栈
的Token序列。本节描述核心设计决策及其Token效率的基本原理。

\subsection{词法设计}

\textbf{Whisper}仅使用ASCII字符集。每个词法Token都是语义意义最大密度单元。
表\ref{tab:lexical}总结了词法元素。

\begin{table}[ht]
\centering
\caption{\textbf{Whisper}的词法元素。}
\label{tab:lexical}
\begin{tabular}{lll}
\toprule
\textbf{类别} & \textbf{Token} & \textbf{示例} \\
\midrule
栈操作 & \texttt{\_ \` @ \$n drop} & \texttt{42 \_ . .} $\rightarrow$ \texttt{42 42} \\
算术 & \texttt{+ - * / \%} & \texttt{3 4 + .} $\rightarrow$ \texttt{7} \\
比较 & \texttt{= < > != <= >=} & \texttt{5 3 > .} $\rightarrow$ \texttt{\#t} \\
逻辑 & \texttt{\& | !} & \texttt{\#t \#f \& .} $\rightarrow$ \texttt{\#f} \\
条件 & \texttt{?? ... | ... ]} & \texttt{\#t ?? 1 | 0 ] .} $\rightarrow$ \texttt{1} \\
循环 & \texttt{\{条件\} \{循环体\} \#} & \texttt{5 \{ \_ 0 > \} \{ 1 - \} \# .} $\rightarrow$ \texttt{0} \\
定义 & \texttt{: 名称 \{ 函数体 \} ;} & \texttt{: sq \{ \_ * \} ;} \\
列表 & \texttt{[ ] @nth @map @fold @each} & \texttt{[1 2 3] \{ \_ * \} @map} \\
字符串 & \texttt{strlen strcat strfind} & \texttt{"ab" "cd" strcat} \\
布尔值 & \texttt{\#t \#f} & \texttt{\#t .} \\
导入 & \texttt{import std/math} & \texttt{import std/list} \\
\bottomrule
\end{tabular}
\end{table}

\subsection{无变量名的参数识别}

一个自然的问题是：\textbf{Whisper}没有变量名，如何区分不同的参数？答案是：
\textbf{栈上的位置就是身份}。Python中，需要写\texttt{a = 3; b = 4}将值绑定
到名称，然后通过名称引用。而在\textbf{Whisper}中，栈位置0（栈顶）\textit{就是}
第一个参数，位置1\textit{就是}第二个参数，以此类推。操作从这些隐式位置消费值。

以计算$(a + b) \times c$为例，其中$a = 5$、$b = 3$、$c = 2$。
Python中需要写\texttt{(a + b) * c}。在\textbf{Whisper}中：

\begin{lstlisting}[language=C,basicstyle=\ttfamily]
5 3 + 2 * .
\end{lstlisting}

栈变化过程：
\texttt{[5]}（压入$a$）$\rightarrow$
\texttt{[5, 3]}（压入$b$）$\rightarrow$
\texttt{[8]}（\texttt{+}弹出两个值，压入$a+b$）$\rightarrow$
\texttt{[8, 2]}（压入$c$）$\rightarrow$
\texttt{[16]}（\texttt{*}弹出两个值，压入结果）$\rightarrow$ 打印。

值\texttt{5}通过其栈位置被识别——它是第一个被压入的值，之后被\texttt{+}
从第二个栈槽中消费。整个过程从未给它赋予名称。

当需要多次使用同一个值时，三个基本栈操作提供基于位置的访问，无需命名：

\begin{itemize}
\item \texttt{\_}（dup，复制）：复制栈顶元素。$a \rightarrow a\ a$。
  \textit{等价于}Python中的\texttt{y = x; z = x}，但仅需一个Token。
\item \texttt{`}（swap，交换）：交换栈顶两个元素。$a\ b \rightarrow b\ a$。
  \textit{等价于}交换两个参数的顺序。
\item \texttt{@}（rot，旋转）：旋转栈顶三个元素。$a\ b\ c \rightarrow b\ c\ a$。
  \textit{等价于}三个临时变量的循环排列。
\item \texttt{drop}（丢弃）：丢弃栈顶元素。
\end{itemize}

以更复杂的$(a + b) \times (a - b)$为例（$a = 5$、$b = 3$），两个值各需使用
两次，却不使用任何变量名：

\begin{lstlisting}[language=C,basicstyle=\ttfamily]
5 3 _ over + ` over - * .
\end{lstlisting}

栈变化如下（栈顶在右）：
\texttt{[5]} $\rightarrow$ \texttt{[5, 3]} $\rightarrow$ \texttt{[5, 3, 3]}
（\texttt{\_}复制$b$）$\rightarrow$ \texttt{[5, 3, 3, 5]}（\texttt{over}复制
$a$）$\rightarrow$ \texttt{[5, 3, 8]}（\texttt{+}）$\rightarrow$ \texttt{[5, 8, 3]}
（\texttt{`}交换）$\rightarrow$ \texttt{[5, 8, 3, 5]}（\texttt{over}）
$\rightarrow$ \texttt{[5, 8, -2]}（\texttt{-}）$\rightarrow$ \texttt{[5, 16]}
（\texttt{*}）。变量$a$和$b$各自从不同栈深度被访问两次，从未被命名。
这种位置式范式消除了所有变量声明和引用Token，在典型Python函数中，
这些Token约占15--20\%。

\subsection{条件语法}

传统语言使用\texttt{if/else}块（Python中每个分支4--6个Token）。
\textbf{Whisper}使用紧凑的三元式语法：

\begin{lstlisting}[language=C,basicstyle=\ttfamily]
条件 ?? 真分支 | 假分支 ]
\end{lstlisting}

对于嵌套条件，额外的\texttt{]}终止符关闭每一层：

\begin{lstlisting}[language=C,basicstyle=\ttfamily]
x _ 0 > ?? "positive"
| x _ 0 < ?? "negative"
| "zero" ] ] .
\end{lstlisting}

这消除了对\texttt{elif}、\texttt{else}、冒号和缩进的需求——
每一个在Python中都要消耗一个Token。

\subsection{函数定义}

函数使用\texttt{: 名称 \{ 函数体 \} ;}定义：

\begin{lstlisting}[language=C,basicstyle=\ttfamily]
: factorial { _ 1 > ?? _ 1 - factorial * | drop 1 ] } ;
\end{lstlisting}

\texttt{\{ \}}分隔符创建引用（延迟执行块），\texttt{: ;}将其绑定到名称。
这比Python的\texttt{def name(args): ... return ...}更紧凑
（每个定义最少节省5个Token开销）。

\subsection{标准库}

我们工作的一个关键洞察是：丰富的标准库可以显著减少Token消耗。
\textbf{Whisper}标准库在三个模块中提供了83个实用词：

\begin{itemize}
\item \texttt{std/math}：\texttt{abs}、\texttt{max}、\texttt{min}、
  \texttt{neg}、\texttt{sq}、\texttt{pow}、\texttt{even?}、\texttt{odd?}、
  \texttt{sqrt}、\texttt{positive?}、\texttt{negative?}、\texttt{zero?}、
  \texttt{inc}、\texttt{dec}、\texttt{double}、\texttt{halve}、
  \texttt{factorial}、\texttt{fib}、\texttt{clamp}、\texttt{between?}
\item \texttt{std/list}：\texttt{sum}、\texttt{prod}、\texttt{first}、
  \texttt{last}、\texttt{tail}、\texttt{rev}、\texttt{mean}、\texttt{sort}、
  \texttt{range}、\texttt{range-to}、\texttt{empty?}、\texttt{contains?}、
  \texttt{max-val}、\texttt{min-val}
\item \texttt{std/str}：\texttt{contains?}、\texttt{rev}、\texttt{repeat}、
  \texttt{capitalize}、\texttt{palindrome?}、\texttt{starts-with?}、
  \texttt{ends-with?}、\texttt{upper}、\texttt{lower}
\end{itemize}

例如，不使用标准库计算绝对值需要8个Token
（\texttt{\_ 0 < ?? 0 swap - ]}），而使用标准库只需1个Token
（\texttt{abs}）。在25个任务的基准测试中，标准库将\textbf{Whisper}代码量
减少了36.6\%（在15个任务的评估集上从273个Token降至173个Token）。

% ═══════════════════════════════════════════════════════════════
\section{安全模型}
\label{sec:security}
% ═══════════════════════════════════════════════════════════════

AI生成代码的一个根本性挑战是信任问题。\textbf{Whisper}通过基于能力的安全模型，
以三层防御体系解决这一问题：

\begin{enumerate}
\item \textbf{编译期}：I/O能力是独立类型（\texttt{Cap(u16)}），
  不可构造、不可强制转换、不可与数据类型混合。从未导入能力模块的程序
  是\textit{可证明的纯函数}。
\item \textbf{加载期}：能力表由宿主环境一次性初始化，运行时不可变。
  在语言层面，动态能力创建是不可能的。
\item \textbf{运行时}：每个能力在受限环境中执行。文件读取能力限制于白名单路径。
  HTTP能力限制于白名单主机。
\end{enumerate}

\textbf{Whisper}包在元数据中声明其所需能力：

\begin{lstlisting}[language=C,basicstyle=\ttfamily]
能力声明: ["@file_read", "@http_post"]
\end{lstlisting}

安装前，用户会看到列出请求能力的授权提示。这类似于Android的权限系统，
但在语言层面而非操作系统层面执行。

\subsection{与现有方法的对比}

现有AI代码安全方法依赖于容器或操作系统层面的沙箱（Docker、gVisor、seccomp）。
这些方法是粗粒度的，作用于整个进程。\textbf{Whisper}的能力模型是细粒度的：
需要读取一个特定文件的程序不能读取任何其他文件，没有能力声明的程序具有零副作用。
这使得\textit{信任但验证}模式成为可能——AI生成的代码可以以精确限定范围的权限执行。

% ═══════════════════════════════════════════════════════════════
\section{实现}
\label{sec:implementation}
% ═══════════════════════════════════════════════════════════════

\subsection{编译器架构}

\textbf{Whisper}工具链包括：

\begin{itemize}
\item \textbf{基于Rust的VM}：拥有40+条操作码的栈式虚拟机，涵盖算术、比较、
  逻辑、控制流、字符串操作、列表操作、JSON解析和浮点数学。
\item \textbf{自举编译器}：\textbf{Whisper}编译器（\texttt{whisperc}）
  本身用\textbf{Whisper}编写（\texttt{lexer.ws}、\texttt{parser.ws}、
  \texttt{classify.ws}）。它将\texttt{.ws}源文件编译为\texttt{.wbin}字节码
  和\texttt{.wasm} WebAssembly。
\item \textbf{多目标代码生成}：编译器支持原生二进制（\texttt{.wbin}）、
  WebAssembly文本（\texttt{.wat}）和WebAssembly二进制（\texttt{.wasm}）
  输出格式。
\end{itemize}

\subsection{二进制格式}

\texttt{.wbin}格式使用单字节操作码和LEB128编码的操作数。
算术运算各占一个字节（\texttt{0x10}--\texttt{0x14}）。
这种紧凑编码进一步减少了生成程序的存储占用。

\subsection{LLM集成}

\textbf{Whisper}被设计为LLM代码生成的输出格式。该语言的最小化语法减少了
解码器的输出长度，直接降低了推理延迟和成本。训练流水线
（第\ref{sec:experiments}节）展示了这一集成在7B参数模型上的效果。

% ═══════════════════════════════════════════════════════════════
\section{实验设置}
\label{sec:experiments}
% ═══════════════════════════════════════════════════════════════

\subsection{研究问题}

我们研究三个问题：

\begin{itemize}
\item \textbf{RQ1}：LLM能否被微调以从自然语言指令生成正确的\textbf{Whisper}代码？
\item \textbf{RQ2}：对于相同任务，\textbf{Whisper}代码是否比等效Python代码
  消耗更少的Token？
\item \textbf{RQ3}：标准库设计对Token效率的影响是什么？
\end{itemize}

\subsection{数据集}

我们构建了包含1,149条指令--响应对的训练数据集，涵盖12个类别：算术、栈操作、
比较与逻辑、条件分支、函数定义、递归、列表操作、字符串操作、控制流、
真实程序、算法问题和语法关键模式。每个示例由自然语言指令
（如``找出两个数中的最小值''）、可选的输入数据和期望的\textbf{Whisper}代码组成。

我们还构建了两个评估基准：

\begin{itemize}
\item \textbf{Eval-15}：15个任务，衡量代码生成准确率，涵盖算术、字符串操作、
  列表操作、递归和算法。
\item \textbf{Benchmark-25}：25个任务，衡量Token效率，涵盖简单表达式
  （``Hello World''）、数学计算、数据结构操作和算法问题。
\end{itemize}

\subsection{模型与训练}

我们使用LoRA（低秩适配）微调了Qwen2.5-Coder-7B-Instruct，配置如下：

\begin{itemize}
\item \textbf{LoRA秩}：$r = 128$，$\alpha = 256$
\item \textbf{量化}：8位（BitsAndBytes）
\item \textbf{训练轮数}：5
\item \textbf{批次大小}：每设备8 $\times$ 2梯度累积 = 16有效批次
\item \textbf{学习率}：$1 \times 10^{-4}$，余弦调度
\item \textbf{最大序列长度}：2,048 Token
\item \textbf{优化器}：Paged AdamW 8位
\item \textbf{GPU}：单卡NVIDIA A6000（48 GB显存）
\end{itemize}

在上述硬件上训练耗时约45分钟。

\subsection{评估指标}

\begin{itemize}
\item \textbf{精确匹配（EM）}：生成的代码与参考\textbf{Whisper}代码在归一化后
  逐字符匹配。
\item \textbf{语法有效性（SV）}：生成的代码具有平衡的\texttt{\{ \}}和
  \texttt{[ ]}分隔符。
\item \textbf{Token数量}：模型分词器对生成代码产生的Token数量。对于
  \textbf{Whisper}，由于语言的纯ASCII单字符操作符特性，每个空格分隔的Token
  通常对应一个模型Token。
\item \textbf{Token减少率}：$(1 - T_{\text{Whisper}} /
  T_{\text{Python}}) \times 100\%$。
\end{itemize}

% ═══════════════════════════════════════════════════════════════
\section{结果}
\label{sec:results}
% ═══════════════════════════════════════════════════════════════

\subsection{代码生成准确率（RQ1）}

表\ref{tab:eval}报告了每个任务的评估结果。微调后的模型实现了40.0\%的精确匹配
准确率（6/15任务）和93.3\%的语法有效性（14/15任务）。

\begin{table}[ht]
\centering
\caption{Eval-15上每个任务的评估结果。}
\label{tab:eval}
\begin{tabular}{llcc}
\toprule
\textbf{任务} & \textbf{描述} & \textbf{EM} & \textbf{SV} \\
\midrule
eval\_01 & 1到20求和 & --- & \checkmark \\
eval\_02 & 两个数的最小值 & \checkmark & \checkmark \\
eval\_03 & 判断正数 & \checkmark & \checkmark \\
eval\_04 & 计算$2^{16}$ & --- & \checkmark \\
eval\_05 & 拼接三个字符串 & \checkmark & \checkmark \\
eval\_06 & 列表最后一个元素 & --- & \checkmark \\
eval\_07 & 列表平均值 & --- & \checkmark \\
eval\_08 & 重复字符串$n$次 & --- & \checkmark \\
eval\_09 & 数字列表转整数 & --- & \checkmark \\
eval\_10 & 字符串包含子串 & \checkmark & \checkmark \\
eval\_11 & 列表元素乘积 & --- & \checkmark \\
eval\_12 & 移除第一个元素 & \checkmark & \checkmark \\
eval\_13 & 两点间距离 & --- & \checkmark \\
eval\_14 & 1到10的平方 & \checkmark & \checkmark \\
eval\_15 & 判断质数 & --- & --- \\
\midrule
\textbf{总计} & & \textbf{6/15} & \textbf{14/15} \\
\bottomrule
\end{tabular}
\end{table}

唯一的语法失败（eval\_15，质数检查）涉及包含多个循环结构的复杂嵌套条件——
这是基准测试中最具挑战性的任务。模型正确识别了所需算法，但在条件嵌套深度上
出现了错误。

\subsection{Token效率（RQ2）}

表\ref{tab:token}展示了Benchmark-25上的Token效率对比。微调后的模型生成的
\textbf{Whisper}代码使用了444个Token，而等效Python代码使用了1,077个——
减少了58.8\%。

\begin{table}[ht]
\centering
\caption{Benchmark-25上的Token数量对比（部分任务）。}
\label{tab:token}
\begin{tabular}{lrrr}
\toprule
\textbf{任务} & \textbf{Whisper} & \textbf{Python} & \textbf{减少率} \\
\midrule
hello\_world & 6 & 10 & 40.0\% \\
factorial & 24 & 58 & 58.6\% \\
fibonacci & 28 & 96 & 70.8\% \\
sum\_list & 19 & 54 & 64.8\% \\
product\_list & 17 & 68 & 75.0\% \\
map\_squares & 18 & 64 & 71.9\% \\
string\_length & 5 & 23 & 78.3\% \\
string\_concat & 8 & 48 & 83.3\% \\
is\_even & 8 & 41 & 80.5\% \\
gcd & 24 & 71 & 66.2\% \\
string\_reverse & 7 & 112 & 93.8\% \\
negate & 6 & 48 & 87.5\% \\
\midrule
\textbf{总计（25个任务）} & \textbf{444} & \textbf{1,077} & \textbf{58.8\%} \\
\bottomrule
\end{tabular}
\end{table}

Token减少在涉及迭代、递归和字符串操作的任务中最为显著——正是
\textbf{Whisper}的栈式后缀表示法和丰富的标准库相比Python更冗长的控制结构
具有最大优势的类别。

\subsection{真实场景对比}

我们设计了五个涵盖不同领域的真实编程任务，对比了\textbf{Whisper}和Python的
Token消耗：

\begin{table}[ht]
\centering
\caption{五个真实场景的Token对比。}
\label{tab:realworld}
\begin{tabular}{lrrr}
\toprule
\textbf{任务} & \textbf{Whisper} & \textbf{Python} & \textbf{结果} \\
\midrule
状态消息格式化 & 6 & 11 & Whisper胜（+45.5\%） \\
JSON配置文件读取 & 21 & 18 & Python胜（$-$16.7\%） \\
API数据抓取与计数 & 36 & 83 & Whisper胜（+56.6\%） \\
词频统计 & 62 & 74 & Whisper胜（+16.2\%） \\
质数筛（100以内） & 43 & 84 & Whisper胜（+48.8\%） \\
\midrule
\textbf{总计} & \textbf{168} & \textbf{270} & \textbf{Whisper胜（+37.8\%）} \\
\bottomrule
\end{tabular}
\end{table}

\textbf{Whisper}在5个任务中的4个上优于Python（\textbf{80\%胜率}），
总体Token减少37.8\%。唯一Python占优的任务（JSON配置文件读取）涉及简单的
键值访问，Python的\texttt{dict[key]}表示法（每次访问1个Token）比
\textbf{Whisper}的\texttt{listfind}模式（每次访问2个Token）更紧凑。

\subsection{标准库的影响（RQ3）}

为了量化标准库的贡献，我们在Eval-15基准上对比了有无扩展标准库的
\textbf{Whisper}代码。表\ref{tab:stdlib}展示了结果。

\begin{table}[ht]
\centering
\caption{标准库扩展前后的Token数量对比（Eval-15）。}
\label{tab:stdlib}
\begin{tabular}{lrrr}
\toprule
\textbf{任务} & \textbf{之前} & \textbf{之后} & \textbf{节省} \\
\midrule
eval\_01（1--20求和） & 26 & 6 & 76.9\% \\
eval\_04（$2^{16}$） & 25 & 6 & 76.0\% \\
eval\_08（重复字符串） & 26 & 6 & 76.9\% \\
eval\_11（乘积） & 10 & 8 & 20.0\% \\
eval\_14（平方列表） & 16 & 11 & 31.2\% \\
\midrule
\textbf{总计（15个任务）} & \textbf{273} & \textbf{173} & \textbf{36.6\%} \\
\bottomrule
\end{tabular}
\end{table}

仅标准库一项就将\textbf{Whisper}代码量减少了36.6\%。最大的节省来自于
将多Token模式替换为单Token标准库调用（例如，
\texttt{\_ 0 < ?? 0 swap - ]} $\rightarrow$ \texttt{abs}，
该模式减少了87.5\%）。

\subsection{成本影响}

按当前API定价（GPT-4o、Claude Opus 4等前沿模型约每百万输出Token \$0.50），
58.8\%的Token减少带来的直接成本节约为：

\begin{itemize}
\item \textbf{每百万次代码生成}：按每次生成平均700个输出Token计算，
  相比Python，\textbf{Whisper}每百万次生成节省约\$206
  （700 $\times$ 58.8\% $\times$ \$0.50/百万Token $\times$ 100万次）。
\item \textbf{企业规模}：对于每天服务1,000万次AI代码补全的组织，
  年度节省超过\$75万
  （1,000万 $\times$ 365 $\times$ 700 $\times$ 58.8\% $\times$ \$0.50/百万Token）。
\end{itemize}

以上估算较为保守；仅考虑输出Token成本，未计算更短的提示词
（\textbf{Whisper}系统提示也更紧凑）或更短解码序列带来的延迟降低等复合效应。

% ═══════════════════════════════════════════════════════════════
\section{相关工作}
\label{sec:related}
% ═══════════════════════════════════════════════════════════════

\subsection{拼接式与栈式语言}

\textbf{Whisper}继承了拼接式语言的丰富传统。Forth~\cite{forth}为资源受限
的嵌入式系统开创了栈式后缀范式。Factor~\cite{factor}通过丰富的标准库和交互式
开发环境使这一方法现代化。Joy~\cite{joy}探索了拼接式编程作为纯函数范式的理论
基础。\textbf{Whisper}与这些前身的不同之处在于其明确针对机器生成而非人类编写
的优化。

\subsection{AI生成代码安全}

近期关于AI生成代码安全的工作主要集中在沙箱~\cite{seccomp,gvisor}和运行时监控
~\cite{code-sandbox}。这些方法在语言外部添加安全层。\textbf{Whisper}的能力
模型将安全集成到类型系统中，使得在合法的\textbf{Whisper}程序中不可能表达
未授权的副作用。

\subsection{Token高效代码生成}

多项研究考察了编程语言语法与LLM代码生成质量之间的关系
~\cite{codex,code-llm-survey}。然而，这些工作将目标语言视为固定的，侧重于
提示工程或模型架构。据我们所知，\textbf{Whisper}是第一门从零开始设计以
最小化LLM Token消耗的语言。

\subsection{LLM优化表示}

并行研究探索了LLM生成代码的二进制或压缩表示~\cite{llm-compression,wasm-generation}。
\textbf{Whisper}的方法与之互补：它提供了一种人类可读（虽然简洁）的文本格式，
可直接映射到紧凑的二进制表示（\texttt{.wbin}）。

% ═══════════════════════════════════════════════════════════════
\section{讨论}
\label{sec:discussion}
% ═══════════════════════════════════════════════════════════════

\subsection{局限性}

\textbf{训练数据规模}：我们1,149个示例的数据集虽然足以进行初步验证，但比主流
语言使用的训练语料库小数个数量级。我们预期随着数据集规模的扩大，精确匹配
准确率将显著提高。

\textbf{模型规模}：我们仅评估了7B参数模型。更大的模型（34B、70B）可能在
\textbf{Whisper}代码生成上实现更高的准确率。

\textbf{可读性权衡}：\textbf{Whisper}为Token优化的语法对人类而言不如Python
或JavaScript可读。我们认为这对于主要供机器消费的机器生成代码来说是可以接受的
权衡。

\textbf{生态系统}：作为一门新语言，\textbf{Whisper}缺乏成熟语言的库生态系统。
能力系统通过支持通过HTTP和文件I/O能力将\textbf{Whisper}程序与现有服务安全
组合，部分缓解了这一问题。

\subsection{未来工作}

\begin{itemize}
\item \textbf{更大规模微调}：将训练数据集扩展至10,000+示例，并在更大的基座
  模型上进行评估。
\item \textbf{多语言对比}：将Token效率基准测试扩展至JavaScript、Java、
  Rust和Go。
\item \textbf{强化学习}：使用执行反馈提升超越精确匹配指标的语义正确性。
\item \textbf{IDE集成}：构建VS Code扩展，提供\textbf{Whisper}语法高亮、
  调试和能力检查。
\item \textbf{生产部署}：在真实代码生成流水线中评估\textbf{Whisper}，
  并进行成本跟踪。
\end{itemize}

% ═══════════════════════════════════════════════════════════════
\section{结论}
\label{sec:conclusion}
% ═══════════════════════════════════════════════════════════════

我们提出了\textbf{Whisper}——一门栈式后缀编程语言，旨在最小化LLM生成代码中
的Token消耗。\textbf{Whisper}通过单字符操作符、紧凑的条件语法和包含83个
实用词的丰富标准库实现Token效率。其基于能力的零信任安全模型提供对I/O的
细粒度控制，支持AI生成代码的安全执行。

我们的实证评估表明，微调后的7B参数模型在\textbf{Whisper}代码生成上实现了
40.0\%的精确匹配准确率和93.3\%的语法有效性，同时生成的代码比等效Python
少消耗58.8\%的Token。在AI推理成本以每Token计价的时代，语言级优化代表了一个
尚未充分开发的手段，可用于降低AI辅助软件开发的经济和环境成本。

% ═══════════════════════════════════════════════════════════════
\bibliographystyle{plain}
\begin{thebibliography}{99}

\bibitem{forth}
L. Brodie, \textit{Starting FORTH}. Prentice-Hall, 1981.

\bibitem{factor}
S. Pestov, D. Ehrenberg, J. Groff, ``Factor: A Dynamic Stack-based
Programming Language,'' \textit{ACM SIGPLAN Notices}, vol. 45, no. 12, 2010.

\bibitem{joy}
M. von Thun, ``Joy: A Functional Language Based on Composition of Functions,''
La Trobe University, Technical Report, 2001.

\bibitem{seccomp}
J. Edge, ``A seccomp overview,'' \textit{LWN.net}, 2015.

\bibitem{gvisor}
Google, ``gVisor: Application Kernel for Containers,''
\textit{USENIX ATC}, 2018.

\bibitem{code-sandbox}
M. Christodorescu et al., ``Towards Secure AI-Generated Code,'' \textit{arXiv
  preprint}, 2024.

\bibitem{codex}
M. Chen et al., ``Evaluating Large Language Models Trained on Code,''
\textit{arXiv:2107.03374}, 2021.

\bibitem{code-llm-survey}
Z. Zheng et al., ``A Survey of Large Language Models for Code,''
\textit{arXiv:2311.08947}, 2023.

\bibitem{llm-compression}
T. Dettmers et al., ``QLoRA: Efficient Finetuning of Quantized LLMs,''
\textit{NeurIPS}, 2023.

\bibitem{wasm-generation}
B. Bichsel et al., ``CodeGen2: Lessons for Enterprise-Grade Code Generation,''
\textit{arXiv:2305.02309}, 2023.

\end{thebibliography}

\end{document}
